\section{ที่มาและความสำคัญ}
\label{sec:intro}

มาตรวัดผลลัพธ์ที่รายงานโดยผู้ป่วย (Patient-Reported Outcome Measures: PROMs) เป็นเครื่องมือมาตรฐานสำหรับวัดสถานะสุขภาพและอาการจากมุมมองของผู้ป่วยโดยตรง และถูกนำไปใช้อย่างแพร่หลายในการติดตามผลการรักษาและการวิจัยทางคลินิก \citep{Black2013} ในสาขาออร์โธปิดิกส์ PROMs ที่ใช้ประเมินความสามารถเชิงหน้าที่ ได้แก่ Disabilities of the Arm, Shoulder and Hand (DASH) และ Toronto Extremity Salvage Score (TESS) ฉบับรยางค์บนและรยางค์ล่าง โดยทุกข้อคำถามมีตัวเลือกคำตอบเป็นมาตราส่วนเชิงอันดับ (ordinal scale) ที่มีคำบรรยายกำกับแต่ละระดับ

การเก็บข้อมูล PROMs มีสองรูปแบบหลัก ได้แก่ การให้ผู้ป่วยตอบแบบสอบถามด้วยตนเองทั้งบนกระดาษและในรูปแบบอิเล็กทรอนิกส์ (electronic PRO: ePRO) และการสัมภาษณ์โดยบุคลากรทางการแพทย์ การตอบแบบสอบถามด้วยตนเองเป็นอุปสรรคต่อผู้สูงอายุ ผู้ที่มีปัญหาด้านการมองเห็นหรือการอ่าน และผู้ป่วยที่มีความบกพร่องของรยางค์บน ซึ่งเป็นกลุ่มเป้าหมายของแบบประเมิน DASH และ TESS โดยตรง ผู้ป่วยกลุ่มนี้จึงต้องได้รับการสัมภาษณ์โดยพยาบาล ซึ่งต้องใช้บุคลากรและเวลาเป็นจำนวนมาก อีกทั้งรูปแบบการเก็บข้อมูลที่แตกต่างกันยังอาจส่งผลต่อคุณภาพของข้อมูล \citep{Bowling2005} ข้อจำกัดเหล่านี้เป็นอุปสรรคต่อการเก็บข้อมูล PROMs อย่างต่อเนื่องในเวชปฏิบัติ

ในช่วงหลายปีที่ผ่านมา เทคโนโลยีการรู้จำเสียงพูดอัตโนมัติ (Automatic Speech Recognition: ASR) \citep{Aung2024, TyphoonWhisperTurbo} และการสังเคราะห์เสียงพูด (Text-to-Speech: TTS) \citep{Pipatanakul2024} ได้รับการพัฒนาให้รองรับภาษาไทยได้ดีขึ้นอย่างต่อเนื่อง ประกอบกับแบบจำลองภาษาขนาดใหญ่ (Large Language Models: LLMs) สามารถประมวลผลภาษาธรรมชาติได้อย่างมีประสิทธิภาพ \citep{Yang2025Qwen3, Thirunavukarasu2023} ความก้าวหน้าดังกล่าวเอื้อต่อการพัฒนาตัวแทนสนทนา (conversational agent) ที่สื่อสารกับผู้ป่วยได้อย่างเป็นธรรมชาติ \citep{Laranjo2018, An2025} งานวิจัยของ Soni และคณะ~\citep{Soni2022} พบว่าผู้เข้าร่วมวิจัยร้อยละ 69.9 เลือกตอบข้อมูลสุขภาพผ่านตัวแทนสนทนามากกว่าแบบฟอร์มออนไลน์ และงานวิจัยของ Yang และคณะ~\citep{Yang2024} แสดงให้เห็นว่าผู้ช่วยเสียง (voice assistant) ที่ทำงานร่วมกับ LLMs สามารถเก็บรวบรวมข้อมูลสุขภาพจากผู้สูงอายุและลดภาระงานของบุคลากรทางการแพทย์ได้ การเก็บข้อมูล PROMs ผ่านการสนทนาด้วยเสียงจึงเป็นแนวทางที่มีศักยภาพทั้งในการลดข้อจำกัดของการตอบแบบสอบถามด้วยตนเองและการลดภาระงานของบุคลากรทางการแพทย์

อย่างไรก็ตาม งานวิจัยส่วนใหญ่ที่ใช้ LLMs เก็บข้อมูลสุขภาพผ่านการสนทนาใช้แบบจำลองที่ให้บริการผ่านระบบคลาวด์ เช่น ChatGPT ซึ่งก่อให้เกิดข้อกังวลด้านความเป็นส่วนตัวของข้อมูล เนื่องจากข้อมูลเสียงและข้อความของผู้ป่วยต้องถูกส่งไปประมวลผลบนเซิร์ฟเวอร์ภายนอก นอกจากนี้ ระบบดังกล่าวยังให้แบบจำลองสร้างข้อความที่ใช้สื่อสารกับผู้ป่วยขึ้นระหว่างการสนทนา ทั้งการสร้างประโยคคำถามจากข้อคำถามในแบบประเมิน \citep{Chen2024} และการสร้างข้อความอธิบายเพิ่มเติมเมื่อคำตอบกำกวม \citep{Guo2025} จึงไม่สามารถควบคุมถ้อยคำที่ผู้ป่วยได้รับให้ตรงตามแบบประเมินต้นฉบับได้ทั้งหมด งานวิจัยนี้จึงมุ่งพัฒนาตัวแทนสนทนาด้วยเสียงภาษาไทยสำหรับเก็บรวบรวมข้อมูล ePRO โดยมีลักษณะสำคัญ 3 ประการ ดังนี้
\begin{enumerate}
    \item ใช้แบบประเมิน DASH และ TESS ฉบับรยางค์บนและรยางค์ล่าง โดยคงถ้อยคำคำถามตามต้นฉบับ เพื่อรักษาความตรงของการวัด (measurement validity)
    \item สื่อสารกับผู้ป่วยด้วยเสียงภาษาไทย โดยใช้ ASR และ TTS ทำงานร่วมกับแบบจำลองภาษา
    \item ประมวลผลทั้งหมดบนเซิร์ฟเวอร์ภายในของโรงพยาบาล เพื่อรักษาความเป็นส่วนตัวของข้อมูลผู้ป่วยและรองรับการใช้งานทางคลินิกได้อย่างปลอดภัย
\end{enumerate}


% แบบประเมิน DASH และ TESS โดยใช้ระบบ ASR และ TTS ทำงานร่วมกับ LLMs ที่ทำงานบนเซิร์ฟเวอร์ภายในของโรงพยาบาล เพื่อให้มั่นใจว่าข้อมูลส่วนบุคคลของผู้ป่วยจะไม่ถูกส่งไปประมวลผลภายนอก 
% และสามารถนำไปใช้งานทางคลินิกได้อย่างปลอดภัย

% งานวิจัยในระยะหลังได้นำแบบจำลองภาษาขนาดใหญ่มาใช้เก็บรวบรวมข้อมูลที่ผู้ใช้รายงานด้วยตนเองผ่านการสนทนา Wei และคณะ~\citep{Wei2024} ใช้แบบจำลองภาษาขนาดใหญ่ขับเคลื่อนตัวแทนสนทนาเพื่อเก็บข้อมูลตามช่องข้อมูล (slot) ที่กำหนด และพบว่าตัวแทนสนทนาเก็บข้อมูลได้ครอบคลุมร้อยละ 79 ของช่องข้อมูลที่ต้องการ Chen และคณะ~\citep{Chen2024} พัฒนาระบบ Chat-ePRO ซึ่งใช้ ChatGPT สร้างบทสนทนาจากแบบประเมิน PROs และสกัดข้อมูลจากคำตอบของผู้ป่วย โดยมีอัตราการตอบกลับสูงกว่าแบบฟอร์มและตัวแทนสนทนาแบบอิงกฎ Guo และคณะ~\citep{Guo2025} พัฒนา HopeBot สำหรับคัดกรองภาวะซึมเศร้าด้วยแบบประเมิน PHQ-9 ซึ่งให้คะแนนสอดคล้องกับการตอบแบบประเมินด้วยตนเองที่ค่า ICC เท่ากับ 0.91 Kaiyrbekov และคณะ~\citep{Kaiyrbekov2025} ใช้ตัวแทนสนทนาทางโทรศัพท์สัมภาษณ์ผู้เข้าร่วมวิจัย แล้วใช้แบบจำลองภาษาขนาดใหญ่สกัดคำตอบจากบทถอดความ ได้ความถูกต้องเฉลี่ยร้อยละ 98 แม้บทถอดความจะมีอัตราความผิดพลาดระดับคำเฉลี่ยร้อยละ 7.7 และ Teferra และคณะ~\citep{Teferra2025} ใช้แบบจำลองภาษาขนาดใหญ่ทำนายคะแนนรายข้อของแบบประเมิน PHQ-8 จากบทสัมภาษณ์ทางคลินิก อย่างไรก็ตาม งานวิจัยเหล่านี้ส่วนใหญ่พัฒนาสำหรับภาษาอังกฤษหรือภาษาจีน ใช้แบบจำลองภาษาขนาดใหญ่ที่ให้บริการผ่านระบบคลาวด์ และให้แบบจำลองสร้างหรือปรับถ้อยคำคำถามระหว่างการสนทนา นอกจากนี้ ยังไม่พบงานวิจัยที่ประยุกต์ใช้ตัวแทนสนทนาด้วยเสียงกับแบบประเมินความสามารถเชิงหน้าที่ทางออร์โธปิดิกส์ และประเมินผลโดยเปรียบเทียบกับการสัมภาษณ์โดยพยาบาล

% การพัฒนาระบบดังกล่าวในบริบทของประเทศไทยจึงมีความท้าทายสำคัญ 4 ประการ ประการแรก คำตอบของผู้ป่วยเป็นภาษาพูดที่มีรูปแบบอิสระ ผู้ป่วยมักบรรยายอาการหรือสถานการณ์ในชีวิตประจำวันแทนการระบุตัวเลือกโดยตรง ระบบจึงต้องจับคู่คำตอบดังกล่าวเข้ากับระดับคะแนนเชิงอันดับของแบบประเมินได้อย่างถูกต้อง ทั้งนี้ ยังไม่มีข้อสรุปว่าแบบจำลองภาษาขนาดใหญ่เหมาะสมกับงานนี้มากกว่าแบบจำลองขนาดเล็กที่ผ่านการปรับจูนหรือไม่ เนื่องจากมีหลักฐานว่าแบบจำลองขนาดเล็กที่ผ่านการปรับจูนให้ผลการจำแนกข้อความดีกว่าแบบจำลองภาษาขนาดใหญ่ที่ไม่ผ่านการปรับจูน (zero-shot) \citep{Bucher2024} ประการที่สอง ความตรงของการวัด (measurement validity) ของแบบประเมินขึ้นอยู่กับถ้อยคำคำถามที่เป็นมาตรฐาน ระบบจึงไม่ควรสร้างหรือดัดแปลงถ้อยคำคำถามขึ้นเอง ซึ่งเป็นข้อจำกัดต่อการใช้แบบจำลองภาษาขนาดใหญ่สร้างบทสนทนาโดยตรงดังเช่นในงานวิจัยที่ผ่านมา ประการที่สาม ข้อมูลสุขภาพเป็นข้อมูลส่วนบุคคลที่มีความอ่อนไหวตามพระราชบัญญัติคุ้มครองข้อมูลส่วนบุคคล พ.ศ. 2562 การส่งข้อมูลเสียงและข้อความของผู้ป่วยไปประมวลผลยังบริการคลาวด์ภายนอกจึงไม่เหมาะสม ระบบต้องประมวลผลด้วยแบบจำลองที่ทำงานบนเซิร์ฟเวอร์ภายใน ซึ่งงานวิจัยของ Wiest และคณะ~\citep{Wiest2024} แสดงให้เห็นว่าแบบจำลองภาษาขนาดใหญ่ที่ทำงานภายในสามารถสกัดข้อมูลทางคลินิกจากข้อความได้อย่างแม่นยำ แต่ต้องดำเนินการภายใต้ทรัพยากรการคำนวณที่จำกัด และประการที่สี่ ยังไม่มีชุดข้อมูลบทสนทนาการประเมิน PROMs ภาษาไทยที่สามารถนำมาใช้พัฒนาและประเมินระบบได้

% วิทยานิพนธ์นี้จึงนำเสนอตัวแทนสนทนาด้วยเสียงภาษาไทยสำหรับเก็บรวบรวมข้อมูล ePRO ตามแบบประเมิน DASH และ TESS สำหรับรยางค์บนและรยางค์ล่าง โดยออกแบบให้แยกการควบคุมบทสนทนาออกจากการตีความคำตอบ แบบประเมินแต่ละฉบับได้รับการนิยามในรูปแบบโครงสร้างข้อมูล JSON ที่ระบุข้อคำถามและระดับคำตอบเชิงอันดับ ส่วนจัดการบทสนทนาแบบเครื่องสถานะจำกัด (finite-state dialogue manager) ควบคุมลำดับการถามให้เป็นไปตามแบบประเมิน และนำเสนอคำถามแต่ละข้อด้วยเสียงที่สร้างไว้ล่วงหน้าพร้อมภาพประกอบ เพื่อให้ถ้อยคำคำถามคงที่ตรงตามต้นฉบับ ระบบแปลงคำตอบของผู้ป่วยเป็นข้อความด้วยการรู้จำเสียงพูดอัตโนมัติ จากนั้นส่วนสกัดตัวเลือกคำตอบ (option extractor) จะจับคู่ข้อความดังกล่าวเข้ากับระดับคะแนนของข้อคำถาม โดยทุกองค์ประกอบทำงานบนเซิร์ฟเวอร์ภายใน ความพร้อมสำหรับการใช้งานทางคลินิกของระบบประเมินจากความสอดคล้องระหว่างคะแนนที่ได้จากระบบกับคะแนนที่ได้จากการสัมภาษณ์โดยพยาบาลในผู้ป่วยจริง ทั้งในระดับรายข้อและระดับคะแนนรวม
